\section{Part F. 선형 탐색(Linear Search)}

% 22
\begin{frame}{탐색 문제를 먼저 명세하자}
  \begin{block}{Search}
    \textbf{Input:} 배열 $A[1..n]$, 찾을 값 $x$\par
    \textbf{Output:} $A[i]=x$인 인덱스 $i$, 없으면 $\NotFound$
  \end{block}
  \begin{itemize}
    \item 중복 원소가 있다면 ``첫 번째 인덱스''처럼 요구사항을 명시한다.
    \item 정렬 여부는 사용할 수 있는 알고리즘을 바꾼다.
  \end{itemize}
\end{frame}

% 23
\begin{frame}[fragile]{선형 탐색(Linear Search)}
  \large
  \begin{algorithmic}[1]
    \Procedure{LinearSearch}{$A[1..n],x$}
      \For{$i \gets 1$ \textbf{to} $n$}
        \If{$A[i]=x$}
          \State \Return $i$
        \EndIf
      \EndFor
      \State \Return $\NotFound$
    \EndProcedure
  \end{algorithmic}
  \vfill
  왼쪽부터 하나씩 비교한다. \keyterm{정렬되지 않은 배열에서도 정확히 동작}한다.
\end{frame}

% 24
\begin{frame}{선형 탐색: $x=12$ 찾기}
  \begin{center}
  \begin{tikzpicture}[scale=1.15,transform shape]
    \foreach \x/\val in {0/3,1/6,2/9,3/12,4/15,5/18,6/21,7/24}{
      \node[cell] (l\x) at (1.16*\x,0) {\val};
    }
    \only<1|handout:0>{\node[active] at (0,0) {3};}
    \only<2|handout:0>{\node[discarded] at (0,0) {3}; \node[active] at (1.16,0) {6};}
    \only<3|handout:0>{\node[discarded] at (0,0) {3}; \node[discarded] at (1.16,0) {6}; \node[active] at (2.32,0) {9};}
    \only<4-|handout:1>{\node[discarded] at (0,0) {3}; \node[discarded] at (1.16,0) {6}; \node[discarded] at (2.32,0) {9}; \node[active] at (3.48,0) {12};}
    \only<1|handout:0>{\node[note,above=5mm of l0] {$i=1$: $3\ne12$};}
    \only<2|handout:0>{\node[note,above=5mm of l1] {$i=2$: $6\ne12$};}
    \only<3|handout:0>{\node[note,above=5mm of l2] {$i=3$: $9\ne12$};}
    \only<4-|handout:1>{\node[font=\small,text=AlgoOrangeText,above=5mm of l3] {$i=4$: 발견!};}
  \end{tikzpicture}
  \end{center}
  \only<1|handout:0>{첫 번째 원소부터 비교한다.}
  \only<2-3|handout:0>{일치하지 않으면 다음 인덱스로 이동한다.}
  \only<4-|handout:1>{찾는 즉시 4를 반환하고 종료한다.}
\end{frame}

% 25
\begin{frame}{선형 탐색 정리}
  \begin{columns}[T,onlytextwidth]
    \column{.48\textwidth}
      \textbf{장점}
      \begin{itemize}
        \item 정렬이 필요 없다.
        \item 한 번의 탐색을 위해 준비 비용이 없다.
        \item 첫 일치에서 즉시 종료한다.
      \end{itemize}
    \column{.48\textwidth}
      \textbf{한계}
      \begin{itemize}
        \item 마지막에 있거나 없으면 전부 본다.
        \item 입력이 두 배면 최악 비교 수도 두 배다.
      \end{itemize}
  \end{columns}
  \begin{takeaway}단순하지만 전제가 적다. ``느리다''가 아니라 \textbf{비교 횟수가 입력 크기에 선형}이라고 말하자.\end{takeaway}
  \muted{정렬되지 않은 배열에서는 하나씩 보는 것이 안전하다. 정렬이라는 추가 정보가 있으면 무엇을 건너뛸 수 있을까?}
\end{frame}
